\section{Classification}
In Classification we learn a mapping $f:\mathbb{R}^{m+1}\mapsto \left\{ 0,1,\ldots,k \right\}$,
drawn from a function class $\mathcal{f}$.

\begin{itemize}
  \item \textbf{Representation}
    \[f \in \mathcal{F}\] 
  \item \textbf{Evaluation}
    \[\text{Loss Measure: }\mathcal{E}\left(f\left(\mathbf{x}\right),y\right)=\mathbb{I}\left[ y\neq f(\mathbf{x}) \right]\] 
    \[\text{Generalisation Loss: }L(\mathcal{E},\mathcal{D}, f)=\mathbb{E}_\mathcal{D}\left[ \mathbb{I}\left[ \mathcal{Y}\neq f(\mathcal{X}) \right] \right]\] 
  \item \textbf{Optimisation}
    \[f^{*}=\argmin_{f \in \mathcal{F}}\left\{ \mathbb{E}_\mathcal{D}\left[ \mathbb{I}\left[ \mathcal{Y}\neq f(\mathcal{X}) \right] \right] \right\}\] 
\end{itemize}
\subsubsection{Classification Approaches}
\begin{figure}[h!]
  \centering
  \includegraphics[width=0.5\textwidth]{classifications}
\end{figure}
\textbf{Distribution-free Classification} paradigm: We seek to learn the
classification boundary directly without probabilistic inference.\ e.g. PAC
approach where we seek to approximate $\mathbb{E}_\mathcal{D}\left[
  \mathbb{I}\left[ \mathcal{Y}\neq f(\mathcal{X}) \right] \right]$.
\textbf{Probabilitistic Classification} paradigm: The classification problem
reduces to an inference problem in which we learn the posterior output class
probability, each $p_\mathcal{Y}\left( y=1|\mathbf{x} \right) $characterising
an inhomogeneous Bernoulli distribution, for which we must parameterise a
distribution for each $\mathbf{x}\in dom(\mathcal{X})$\\
\textbf{Generative Classification}: To learn $p_\mathcal{Y}(y|\mathbf{x})$
indirectly by inferring the likelihood $p_\mathcal{X}(\mathbf{x}|y)$ and the
prior $p_\mathcal{Y}(y)$ for each class separately. Requires the most number of
parameters to learn, but allows us to learn $p_\mathcal{X}(\mathbf{x})$.\\
\textbf{Discriminative Classification}: To learn $p_\mathcal{Y}(y|\mathbf{x})$
directly. Less demanding in terms of number of parameters to learn, but is
inflexible, requiring re-running the lagorithm for changes of the loss
function and does not deal well with class imbalances.
\subsection{Baye's Optimal Classifier}
For the binary classification problem, we have Baye's Optimal Classifier
minimising misclassification loss: 
\[f*(\mathbf{x})=
\begin{cases}
  1&\text{if}\quad p_\mathcal{Y}(y=1|\mathbf{x})\ge 0.5\\
  0&\text{if}\quad p_\mathcal{Y}(y=1|\mathbf{x})< 0.5
\end{cases}\] 
\begin{proof}
  \begin{align*}
    f^{*}=&\argmin_{f \in \mathcal{F}}\left\{ \mathbb{E}_\mathcal{D}\left[ \mathbb{I}\left[ \mathcal{Y}\neq f(\mathcal{X}) \right] \right] \right\}\\
         =&\argmin_{f \in \mathcal{F}}\left\{ \sum_{y\in \left\{ 0,1 \right\}} \int p_\mathcal{Y}(y|x)p_\mathcal{X}(\mathbf{x})\mathbb{I}\left[ y\neq f(\mathbf{x}) \right] \,d\mathbf{x} \right\}\\
         =&\argmin_{f \in \mathcal{F}}\left\{ \sum_{y\in \left\{ 0,1 \right\}} \int p_\mathcal{Y}(y|x)p_\mathcal{X}(\mathbf{x})\left( f\left(\mathbf{x}\right) \left(1-y\right) +\left(1-f\left(\mathbf{x}\right)\right)y\right)  \,d\mathbf{x} \right\}\\
         =&\argmin_{f \in \mathcal{F}}\left\{ \int \left( p_\mathcal{Y}\left( y=1|\mathbf{x} \right) \left(1-f\left(\mathbf{x}\right)\right) +p_\mathcal{Y}\left( y=0|\mathbf{x} \right) f\left( \mathbf{x} \right)  \right) p_\mathcal{X}(\mathbf{x})\,d\mathbf{x}  \right\}\\
         =&\argmin_{f \in \mathcal{F}}\left\{ \int \left( p_\mathcal{Y}\left( y=1|\mathbf{x} \right) \left(1-f\left(\mathbf{x}\right)\right) +\left( 1-p_\mathcal{Y}\left( y=1|\mathbf{x} \right) \right)  f\left( \mathbf{x} \right)  \right) p_\mathcal{X}(\mathbf{x}) \,d\mathbf{x}  \right\}\\
         =&\argmin_{f \in \mathcal{F}}\left\{ \int \left( p_\mathcal{Y}\left( y=1|\mathbf{x} \right) +f\left( \mathbf{x} \right) \left( 1-2p_\mathcal{Y}\left( y=1|\mathbf{x} \right) \right) \right) p_\mathcal{X}(\mathbf{x}) \,d\mathbf{x}  \right\}
  \end{align*}
  Let $M(\mathbf{x})=p_\mathcal{Y}\left( y=1|\mathbf{x} \right) +f\left( \mathbf{x} \right) \left( 1-2p_\mathcal{Y}\left( y=1|\mathbf{x} \right) \right)$.\\
  We need to find function $f^{*}$ that is optimal for all $\mathbf{x}\in dom(\mathcal{X})$, 
  \[f^{*}(\mathbf{x})=\argmin_{f(\mathbf{x})}\left\{ M(\mathbf{x}) \right\}\]
  If $f^{*}(\mathbf{x})=1$ optimality implies:
  \begin{align*}
    1-p_\mathcal{Y}\left( y=1|\mathbf{x} \right) \le p_\mathcal{Y}\left( y=1|x \right) 
    p_\mathcal{Y}\left( y=1|\mathbf{x} \right) \ge 0.5
  \end{align*}\qed%
\end{proof}
\begin{remark}
  WLOG, we may consider the Bayes Optimal Classifier for Generative Classification as follows:
  \begin{align*}
    f^{*}(\mathbf{x})=&\argmax_{y\in \left\{ 0,1 \right\}}\left\{ p_\mathcal{Y}(y|\mathbf{x}) \right\}\\
    =&\argmax_{y\in \left\{ 0,1 \right\}}\left\{ \frac{p_\mathcal{X}\left( \mathbf{x}|y \right)P_\mathcal{Y}(y) }{\sum_{y\in \left\{ 0,1 \right\}} p_\mathcal{X}(\mathbf{X}|y)p_\mathcal{Y}(y)} \right\}\\
    =&\argmax_{y\in \left\{ 0,1 \right\}}\left\{ p_\mathcal{X}(\mathbf{x}|y)p_\mathcal{Y}(y) \right\}
  \end{align*}
\end{remark}
\subsection{Logistic Regression}
We seek to model $p_\mathcal{Y}(y=1|\mathbf{x})$ with a function
$f:\mathbb{R}^{m}\mapsto [0,1]$. In Logistic Regression, we assume that the probability can be
modelled by the logistic sigmoid:
\[S(x)=\frac{1}{1+e^{-\mathbf{w\cdot x}}}.\]
\begin{figure}[h!]
  \centering
  \includegraphics[width=0.6\textwidth]{logistic-sigmoid}
\end{figure}
\begin{itemize}
  \item \textbf{Representation}
    \[\mathcal{F}=\left\{ f_\mathbf{w}\left( \mathbf{x} \right) =\mathbb{I}\left[ p_\mathcal{Y}\left( y=1|\mathbf{x} \right) \ge 0.5 \right] | p_\mathcal{Y}\left( y=1|\mathbf{x} \right) =\frac{1}{1+e^{-\mathbf{w\cdot x}}},\mathbf{w}\in \mathbb{R}^{m+1} \right\}\] 
  \item \textbf{Evaluation}
      \[\mathcal{L}(\mathbf{w})= \sum_{i=1}^{n} y^{(i)}\mathbf{w\cdot
      x^{(i)}}-\ln\left( 1+e^{\mathbf{w\cdot x^{(i)}}} \right) \]
  \item \textbf{Optimisation}
    \[\mathbf{w}_{MLE}=\argmin_{\mathbf{w}}\left\{ \sum_{i=1}^{n} \ln\left( 1+e^{\mathbf{w\cdot x}^{(i)}} \right) - y^{(i)}\mathbf{w\cdot x}^{(i)}\right\}\] 
\end{itemize}
Rearranging, we have
\begin{align*}
  p_\mathcal{Y}\left( y=1|\mathbf{x} \right) =&\frac{1}{1+e^{-\mathbf{w\cdot x}}}\\
  e^{-\mathbf{w\cdot x}}=&\frac{1}{p_\mathcal{Y}(y=1|\mathbf{x})}-1\\
  \mathbf{w\cdot x}=&\ln\left( \frac{p_\mathcal{Y}(y=1|\mathbf{x})}{1-p_\mathcal{Y}(y=1|\mathbf{x})} \right)\\
  =&logit(p_\mathcal{Y}(y=1|\mathbf{x}))
\end{align*}
\begin{remark}
$\mathbf{w\cdot x}=0$ defines a linear discriminant separating hyperplane.
\end{remark}
\subsubsection{Evaluation}
\begin{multicols}{2}
  [With the logistic regression model, we have the following:]
   \begin{align*}
    \mathcal{Y}\left( y=1|\mathbf{x} \right)=\frac{1}{1+e^{-\mathbf{w\cdot x}}}\\\\
   \end{align*}
   \begin{align*}
     p_\mathcal{Y}\left( y=0|\mathbf{x} \right)=&\frac{e^{-\mathbf{w\cdot x}}}{1+e^{-\mathbf{w\cdot x}}}\\
     =&\frac{1}{1+e^{\mathbf{w\cdot x}}}
   \end{align*}
\end{multicols}
We can then derive the Cross-Entropy loss function for logistic regression:
\begin{align*}
\mathcal{L}(\mathbf{w})=&\ln\left( \prod_{i=1}^{n} p_\mathcal{Y}\left( y^{(i)}|\mathbf{x}^{(i)} \right)   \right) \\
=&\sum_{i=1}^{n} \ln\left( p_\mathcal{Y}\left( y^{(i)}|\mathbf{x}^{(i)} \right)  \right) \\
=&\sum_{i=1}^{n} y^{(i)}\ln\left( p_\mathcal{Y}\left( y^{(i)}=1|\mathbf{x}^{(i)} \right)  \right) + \left( 1-y^{(i)}\right) \ln\left( p_\mathcal{Y}\left( y^{(i)}=0|\mathbf{x}^{(i)} \right)  \right)  \\
=&\sum_{i=1}^{n} \ln\left( p_\mathcal{Y}\left( y^{(i)}=0|\mathbf{x}^{(i)} \right)  \right) + y^{(i)} \ln\left( \frac{p_\mathcal{Y}\left( y^{(i)}=1|\mathbf{x}^{(i)} \right)}{p_\mathcal{Y}\left( y^{(i)}=0|\mathbf{x}^{(i)} \right)} \right)  \\
=&\sum_{i=1}^{n} y^{(i)}\mathbf{w\cdot x^{(i)}}-\ln\left( 1+e^{\mathbf{w\cdot x^{(i)}}} \right) 
\end{align*}
Hence we seek $\mathbf{w}_{MLE}$:
\[\mathbf{w}_{MLE}=\argmin_{\mathbf{w}}\left\{ \sum_{i=1}^{n} \ln\left( 1+e^{\mathbf{w\cdot x}^{(i)}} \right) - y^{(i)}\mathbf{w\cdot x}^{(i)}\right\}\] 
\subsubsection{Motivation}
We assume some latent random variable $\mathcal{Y}^{*}$ such that
$y^{*}=\mathbf{w\cdot x}+\epsilon$ where $\epsilon\sim Logistic(0,1)$. Then
\[\varepsilon\sim Logistic(\mu,s)\quad\text{where:}\quad\mu\in \mathbb{R},s>0.\]
\begin{definition}[Logistic Distribution]
  The Logistic Distribution has the characteristic pdf:
\[f_\epsilon\left( \varepsilon;\mu,s \right) =\frac{e^{-\frac{\varepsilon-\mu}{s}}}{s{\left( 1+e^{-\frac{\varepsilon-\mu}{s}} \right) }^{2}}.\] 
\[E_\mathcal{D}[\epsilon]=\mu.\] 
\[P(\epsilon<z)=\frac{1}{1+e^{-\frac{z-\mu}{s}}}.\] 
\end{definition}
We link a latent variable outcome to an output assuming the classification
model:
\[y^{(i)}=\begin{cases}
  1&\text{if }\quad y^{*(i)}\ge 0\\
  0&\text{if }\quad y^{*(i)}< 0\\
\end{cases}.\]
\begin{align*}
  p_\mathcal{Y}\left( y^{(i)}=1|\mathbf{x}^{(i)} \right)=&P\left(\mathcal{Y}^{*(i)}\ge 0|\mathbf{x}^{(i)}\right)\\
  =&P\left( \mathbf{w\cdot x}^{(i)}+\varepsilon^{(i)} \ge 0\right)\\
  =&1-\frac{1}{1+e^{\mathbf{w\cdot x}^{(i)}}}\\
  =&\frac{e^{\mathbf{w\cdot x}^{(i)}}}{1+e^{\mathbf{w\cdot x}^{(i)}}}\\
  =&\frac{1}{1+e^{-\mathbf{w\cdot x}^{(i)}}}\\
\end{align*}
\subsubsection{Optimisation}
However, the $\mathbf{w}_{MLE}$ we seek has no analytic solution. We may apply
a numerical technique to optimise. This optimisation is a convex problem, hence the
local optimum found would also be a global optimum.
\[\nabla_\mathbf{w}^{2} \mathcal{L}\ge 0.\] 
\begin{proof}
  \begin{align*}
    \nabla_\mathbf{w} (-y\mathbf{w\cdot x})=&-y\mathbf{x}\\
    \nabla_\mathbf{w}^{2} (-y\mathbf{w\cdot x})=&\mathbf{0}\\
    \implies \mathcal{H}=&\mathbf{0}
  \end{align*}
  $-y\mathbf{w\cdot x}$ is convex.\\
  \begin{align*}
  \nabla_\mathbf{w}\left(  \ln(1+e^{\mathbf{w\cdot x}}) \right)=&\frac{e^{\mathbf{w\cdot x}}}{1+e^{\mathbf{w\cdot x}}} \mathbf{x}\\
  \frac{\partial \ln(1+e^{\mathbf{w\cdot x}})}{\partial w_i} =&\frac{e^{\mathbf{w\cdot x}}}{1+e^{\mathbf{w\cdot x}}} x_i\\
  \frac{\partial^{2} \ln(1+e^{\mathbf{w\cdot x}})}{\partial w_i\partial w_j} =&\frac{e^{\mathbf{w\cdot x}}}{{(1+e^{\mathbf{w\cdot x}})}^{2}} x_i x_j\\
  \mathcal{H}=&\frac{e^{\mathbf{w\cdot x}}}{{(1+e^{\mathbf{w\cdot x}})}^{2}}\mathbf{xx}^{T}\\
  \implies\mathcal{H}\succeq&0
  \end{align*}
  The objective is convex.\qed%
\end{proof}
\subsubsection{Linearly Separable Data}
Consider that the training data is linearly separable. Consider the objective
function for some $w=c\tilde{w}, c>0$:
\[\sum_{i=1}^{n} \ln\left( 1+e^{c\mathbf{\tilde{w}\cdot x}^{(i)}} \right)-y^{(i)}c\mathbf{\tilde{w}\cdot x}^{(i)}\] 
Taking the derivative with respect to $c$, we have:
 \[\sum_{i=1}^{n} \mathbf{\tilde{w}\cdot x}^{(i)}\left( \frac{e^{c\mathbf{\tilde{w}\cdot x}^{(i)}}}{1+e^{c\mathbf{\tilde{w}\cdot x}^{(i)}}}-y^{(i)} \right).\]
Since the data is well classified, there always exists a $\tilde{w}$ such that
the objective has a negative derivative with respect to $c$, and gradient
descent will cause $c$ to grow without bound. The sigmoid function approaches
a Heaviside function. This is a case of overfitting and can be circumvented
through regularisation:
\[\sum_{i=1}^{n} \ln\left( 1+e^{c\mathbf{\tilde{w}\cdot x}^{(i)}} \right)-y^{(i)}c\mathbf{\tilde{w}\cdot x}^{(i)}+\lambda c^{2}\left\lVert \mathbf{\tilde{w}} \right\rVert_2^{2} \]
We may discern a unique stationary point with respect to $c$ with the
derivative, using the addition of a Ridge Regulariser. Apart from the linearly
separable non-regularised case, Logistic Regression in general leads to unique
solutions.
\subsubsection{Multinomial Logistic Regression}
Supposing we have $k$ classes such that $y\in \left\{ 1,\ldots,k \right\}$, the Bayes
Optimal Classifier becomes:
\[f^{*}(\mathbf{x})=\argmax_{y\in \left\{ 1,\ldots,k \right\}}\left\{ p_\mathcal{Y}(y|\mathbf{x}) \right\}.\] 
The model is now defined using the softmax function, and we seek to learn an
inhomogeneous multinomial distribution:
\[p_\mathcal{Y}(y=k|\mathbf{x})=\frac{e^{\mathbf{w}_j\cdot \mathbf{x}}}{\sum_{j=1}^{k} e^{\mathbf{w}_j\cdot\mathbf{x}}}.\] 
The maximum likelihood solution will now make use of the multinomial cross
entropy which results in a convex optimisation problem.

\subsection{Naïve Bayes}
In Naïve Bayes, we simplify the likelihood by assuming the conditional
independence assumption of the input attributes given $\mathcal{Y}$:
\[p_\mathcal{X}(\mathbf{x}|y)=\prod_{i=1}^{m} p_\mathcal{X_i} (x_i|y)\] 
Supposing $\mathcal{Y}$ is a Bernoulli random variable,
$y\sim Bern(\theta_y)$. We find the log-likelihood function:
\begin{align*}
\mathcal{L}\left( \theta \right) =& \ln\left( \prod_{i=1}^{n} p_\mathcal{Y}\left( y^{(i)};\theta_y \right) \right)\\
=&\sum_{i=1}^{n} \ln\left( p_\mathcal{Y}\left( y^{(i)};\theta_y \right)  \right)\\
=&\sum_{i=1}^{n} y^{(i)}\ln\theta_y+\left(1-y^{(i)}\right)\ln\left( 1-\theta_y \right) 
\end{align*}
Hence we seek $\theta_{yMLE}$ such that:
\[\theta_{yMLE}=\argmax_{\theta_y}\left\{ \sum_{i=1}^{n} y^{(i)}\ln\theta_y+\left(1-y^{(i)}\right)\ln\left( 1-\theta_y \right) \right\}\] 
Taking the derivative and setting to 0, we have:
\begin{align*}
\nabla_{\theta_y}L=&\sum_{i=1}^{n} \frac{y^{(i)}}{\theta_y}-\frac{(1-y^{(i)})}{1-\theta_y}\\
\implies\theta_{yMLE}=&\frac{n_1}{n}
\end{align*}
\subsubsection{Categorical Naïve Bayes}
Suppose that the features $x_i$ are discrete-valued and take $m_i$ different
values, such that each outcome $x_i$ is taken from the set ${\left\{ x_{ij}
\right\}}_{j=1}^{m_i}$
\begin{itemize}
  \item \textbf{Representation}
      \[\mathcal{F}=\left\{ f_{\theta_y,\left\{ \theta_{ijk} \right\}} (\mathbf{x})=\argmax_{y\in \left\{ 0,1 \right\}} \left\{ p_\mathcal{Y}(y)\prod_{i=1}^{m} p_\mathcal{X_i} (x_i|y) \right\} \mid p_\mathcal{Y}(y=1)=\theta_y,{\left\{ p_\mathcal{X_i}(x_{ij}|k)=\theta_{ijk} \right\}}^{m,m_i,1}_{i=1,j=1,k=0} \right\}.\] 
  \item \textbf{Evaluation}
    \[\ln\left( L(\theta_y) \right)\quad\text{and}\quad {\left\{ \ln\left( L\left( \theta_{ik} \right)  \right) \right\}}^{m,1}_{i=1,k=0}\] 
  \item \textbf{Optimisation}
    \[\theta_{yMLE}=\frac{n_1}{n}\quad\text{and}\quad {\left\{\theta_{ijkMLE}=\frac{n_{ijk}}{n_k}\right\}}^{m,m_i,1}_{i=1,j=1,k=0}\] 
\end{itemize}
Each $(x_i|y=k)$ is a categorical random variable, of which we seek to
parameterise its categorical distribution.
\begin{definition}[Categorical Distribution]
  The Categorical Distribution is a generalisation of the Bernoulli
  Distribution with more than 2 discrete outcomes ($m_i$ discrete outcomes):
  \[(\mathbf{x}_i|y=k)\sim Categorical(\Theta_{ik})\]
  \[\Theta_{ik}\text{ has elements }{\left\{ \theta_{ijk} \right\}}_{j=1}^{m_i}\]
  \[p_\mathcal{X_i}\left( \mathcal{X}_i=x_{ij}|y=k;\Theta_{ik} \right)=\theta_{ijk} \] 
  \[\sum_{j=1}^{m_i} \theta_{ijk}=1\]
\end{definition}
By forming the log-likelihood, and performing a constrained optimisation of the
resulting function, we can obtain:
\[\theta_{ijkMLE}=\frac{n_{ijk}}{n_k}.\]
where $n_{ijk}$ is the number of training points for which the $i$-th attribute
belongs to the $j$-th category and $\mathcal{Y}=k$, and $n_k$ is the number of
training points of which $\mathcal{Y}=k$.
\subsubsection{Additive Smoothing}
To remedy the problem such that any category does not contain any instances,
and prevent a case of overfitting where we estimate the probability of the
event to be zero, we may utilise additive smoothing, adjusting the MLE
estimates such that:
\[\theta_{ijk}=\frac{n_{ijk}+\alpha}{n_k+\alpha J}.\] 
\[\theta_y=\frac{n_1+\alpha}{n+\alpha K}.\]
where:
\begin{itemize}
  \item $J=$ \# of distinct values outcomes of $\mathcal{X}_i$ can take ($m_i$ in this case)
  \item $K=$ \# of distinct values outcomes of $\mathcal{Y}$ can take
  \item $\alpha=$ the strength of smoothing
\end{itemize}
\subsubsection{Gaussian Naïve Bayes}
Suppose that the features $x_i\in \mathbb{R}$ are continuous valued, such that
each $(\mathcal{X}_i|y=k)$ is a Gaussian random variable, $(x_i|y=k)\sim
\mathcal{N}\left( \mu_{ik},\sigma_{ik} \right) $. We then find the
log-likelihood:
\begin{align*}
  \mathcal{L}(\theta)=&\ln\left( \prod_{j=1}^{n_k} \frac{1}{\sqrt{2\pi\sigma_{ik}^{2}}}\exp\left(-\frac{{(x_i^{(j)}-\mu_{ik})}^{2}}{2\sigma_{ik}^{2}}\right)  \right) \\
  =&-n_k \ln\sigma_{ik}-\sum_{j=1}^{n_k} \left(-\frac{{(x_i^{(j)}-\mu_{ik})}^{2}}{2\sigma_{ik}^{2}}\right)+const.
\end{align*}
We can optimise the log-likelihood with:
\[\mu_{ikMLE}=\sum_{j=1}^{n_k} \frac{x_i^{(j)}}{n_k}.\] 
\[\sigma_{ikMLE}^{2}=\frac{1}{n_k}\sum_{j=1}^{n_k} {(x_i^{(j)}-\mu_{ik})}^{2}.\] 
\begin{itemize}
  \item \textbf{Representation}
    \[\mathcal{F}=\left\{ f_{\theta_y,\left\{ \theta_{ijk} \right\}} (\mathbf{x})=\argmax_{y\in \left\{ 0,1 \right\}} \left\{ p_\mathcal{Y}(y)\prod_{i=1}^{m} \mathcal{N}\left(X_i;\mu_{ik},\sigma_{ik}\right) \right\} \mid p_\mathcal{Y}(y=1)=\theta_y,{\left\{ \mu_{ik}\in \mathbb{R},\sigma_{ik}>0 \right\}}^{m,1}_{i=1,k=0} \right\}.\] 
  \item \textbf{Evaluation}
    \[\ln\left( L(\theta_y) \right)\quad\text{and}\quad {\left\{ \ln\left( L\left( \mu_{ik},\sigma_{ik}\right)  \right) \right\}}^{m,1}_{i=1,k=0}\] 
  \item \textbf{Optimisation}
    \[\theta_{yMLE}=\frac{n_1}{n}\quad\text{and}\quad {\left\{\mu_{ikMLE}=\sum_{j=1}^{n_k} \frac{x_i^{(j)}}{n_k},\sigma_{ikMLE}=\frac{1}{n_k}\sum_{j=1}^{n_k} {(x_i^{(j)}-\mu_{ik})}^{2} \right\}}^{m,1}_{i=1,k=0}\] 
\end{itemize}
\begin{remark}
  GNB results in a similar form of linear discriminant as LR\@. However, they do
  not result in the same classifier, as GNB makes certain model assumptions
  that defines a dependence between $w_0$ and $w_i$.
\end{remark}

\subsection{Comparisons}
Using a generative classifier requires learning more parameters than
discriminative classifiers, but the generative approach is often intractable
and require certain model assumptions (e.g. Conditional Independence). Since LR
makes fewer model assumptions than a generative one such as NB, LR is
considered more robust and less sensitive to modelling choices than NB, but NB
requires less data for similar levels of convergence.
